\documentclass[aspectratio=169,10pt]{beamer}
\usepackage[brazil]{babel}

\usetheme{UELdefesadark}

\title[Título curto]{Título do trabalho de conclusão de curso}
\subtitle{Subtítulo, se houver}
\author[Nome do Autor]{Nome do Autor}
\orientador{Prof. Dr. Nome do Orientador}
\banca{Prof. Dr. Nome do Membro 2\\Profa. Dra. Nome do Membro 3}
\curso{Ciência da Computação}
\date{\today}

\begin{document}

\maketitle

\begin{frame}{Roteiro}
  \tableofcontents
\end{frame}

\section{Introdução}

\begin{frame}{Motivação}{O problema e por que importa}
  \begin{itemize}
    \item Contexto do problema
    \item Lacuna identificada na literatura
    \item Contribuição deste trabalho
  \end{itemize}
\end{frame}

\begin{frame}{Objetivos}
  \begin{block}{Objetivo geral}
    Uma frase que resume o que o trabalho pretende alcançar.
  \end{block}
  \begin{enumerate}
    \item Objetivo específico um
    \item Objetivo específico dois
  \end{enumerate}
\end{frame}

\section{Método}

\begin{frame}{Abordagem}
  \begin{itemize}
    \item Etapa 1: \alert{coleta} dos dados
    \item Etapa 2: implementação
    \item Etapa 3: avaliação
  \end{itemize}
\end{frame}

\section{Resultados}

\begin{frame}{Resultados}
  \centering
  \begin{tabular}{lrr}
    \hline
    Método & Precisão & Tempo (s) \\ \hline
    Linha de base & 0,81 & 12,4 \\
    Proposto & 0,89 & 9,7 \\ \hline
  \end{tabular}
  \par\medskip
  \begin{exampleblock}{Destaque}
    O método proposto melhora a precisão em 8 pontos percentuais.
  \end{exampleblock}
\end{frame}

\section{Conclusão}

\begin{frame}{Conclusões e trabalhos futuros}
  \begin{itemize}
    \item O que foi alcançado
    \item Limitações
    \item Próximos passos
  \end{itemize}
\end{frame}

\obrigado[nome.autor@exemplo.uel.br]

\end{document}
